\documentclass[main]{subfiles}
\usepackage[sols]{workbook}
\numbering{ws2-11}

\begin{document}

\ifSubfilesClassLoaded{%
  \chaptermark{\chaptertwo}
}{}

\section{Power Series} \label{ws2-11}
\begin{framed}
\centerline{\textbf{Lesson Goals}}

You should be able to:
\begin{itemize}
\item Determine whether a given expression is a power series
\item Understand the possible types of convergence behavior for a power series
\item Use the Ratio Test and other series convergence tests to find the radius of convergence, largest open interval of convergence, and the interval of convergence for a given power series
\item Understand the relevance and importance of the interval of convergence of a power series
\item Understand what it means to be a power series representation of a function, and determine the values of $x$ for which a certain power series representation is valid
\end{itemize}
\end{framed}    

\begin{enumerate}
   
  \item  \note{This problem is here to remind students that they have
    already found the interval of convergence for many power
    series, even though we haven't used that terminology yet.}


In previous classes and PSets, we used the Ratio Test along with other strategies to determine the convergence for the Taylor series for $\ln(1+x)$, $\displaystyle\frac{1}{1-x}$, $e^x$, $\cos(x)$, and $\sin(x)$ centered at 0. For each Taylor series, state the values of $x$ for which the series converges. %Write your answers in interval notation.

\begin{enumerate}
\item 
\begin{problem}[\vfill]
The Taylor series for $\ln(1 + x)$ centered at 0: $\displaystyle \sum_{k=1}^\infty
   \frac{(-1)^{k+1}}{k}x^k = x - \frac{1}{2}x^2 + \frac{1}{3}x^3 - \cdots$
\end{problem}

\begin{solution}
\fbox{$(-1,1]$}
\end{solution}

\item \begin{problem}[\vfill]
   The Taylor series for $\displaystyle \frac{1}{1 - x}$ centered at 0:
    $\displaystyle \sum_{k=0}^\infty x^k = 1 + x + x^2 + \cdots$
\end{problem}

\begin{solution}
The Taylor series converges to $\frac{1}{1-x}$ for $x$ in $(-1,1)$, and diverges otherwise.
\end{solution}

\item 
\begin{problem}[\vfill]
The Taylor series for $e^x$ centered at 0: $\displaystyle \sum_{k=0}^\infty
    \frac{x^k}{k!} = 1+x + \frac{1}{2!}x^2 + \frac{1}{3!}x^3 + \cdots$
\end{problem}


\begin{solution}
\fbox{$(-\infty,\infty)$}
\end{solution}


\item 
\begin{problem}[\vfill]
The Taylor series for $\cos(x)$ centered at 0: $\displaystyle \sum_{k=0}^\infty
    \frac{(-1)^k}{(2k)!}x^{2k} = 1- \frac{1}{2!}x^2 + \frac{1}{4!}x^4 - \cdots$
\end{problem}


\begin{solution}
\fbox{$(-\infty,\infty)$}
\end{solution}



\item 
\begin{problem}[\vfill]
The Taylor series for $\sin(x)$ centered at 0: $\displaystyle \sum_{k=0}^\infty
    \frac{(-1)^k}{(2k+1)!}x^{2k+1} = x- \frac{1}{3!}x^3 + \frac{1}{5!}x^5 - \cdots$
\end{problem}

\begin{solution}
\fbox{$(-\infty,\infty)$}
\end{solution}

\end{enumerate}
  \end{enumerate}
  
  
  
  
 \pspace{\newpage}
  
  \begin{problemonly}
   \begin{framed}
\textbf{Definition:}  A \uline{power series} centered at the number $c$ is a series of the form
  $$\displaystyle \boxed{\sum_{k=0}^\infty a_k (x - c)^k=a_0 + a_1 (x - c) +
  a_2 (x - c)^2 + \cdots }$$ where $x$ is a variable and the $a_k$ are
  constants. We can informally think of a power series as an ``infinite polynomial".
  
  Note: The Taylor series we have seen are all particular examples of power series. 
  \end{framed}
  \end{problemonly}

  \begin{enumerate}[resume]
  \item 
  \note{This is just to help students understand the definition of power
  series.} 
  
  Which of the following are power series? (\textbf{Hint: As always, write out the terms in $\cdots$ notation!})
  \probonly{}
  \begin{enumerate}
    \item

    \note{Students tend to get pretty confused with the notation here;
    they look at the definition and focus on the wrong things, like
    the fact that, in the definition, the sum starts with $k = 0$, and
    this series doesn't.  Similarly, the $1 + 2k$ exponent instead of
    just $k$ misleads them into thinking that it's not a power series.
    I emphasize that writing the series in $\cdots$ notation is the
    best way to get a grasp on what it is.}

    \begin{problem}[0.2in]
    $\displaystyle \sum_{k=1}^\infty \sqrt{k} (x - 3)^{1 + 2k}$
    \end{problem}

    \begin{solution}
    This is a power series centered at 3.  It is more evident if we
    write it out in $\cdots$ notation: $\displaystyle
    \sum_{k=1}^\infty \sqrt{k} (x - 3)^{1 + 2k} = (x - 3)^3 +
    \sqrt{2} (x - 3)^5 + \sqrt{3} (x - 3)^7 + \cdots$.
    \end{solution}

    \item
    \begin{problem}[0.2in]
    $\displaystyle \sum_{k=1}^\infty \sqrt{k} (x - 3)^{1 - 2k}$
    \end{problem}

    \begin{solution}
    $\displaystyle \sum_{k=1}^\infty \sqrt{k} (x - 3)^{1 - 2k} = (x -
    3)^{-1} + \sqrt{2} (x - 3)^{-3} + \sqrt{3} (x - 3)^{-5} + \cdots$,
    which is not a power series because it has terms with $(x - 3)$
    raised to negative powers.
    \end{solution}


    \item
    \begin{problem}[0.2in]
    $\displaystyle \sum_{k=0}^\infty \sqrt{x} (x - 5)^{k}$
    \end{problem}

    \begin{solution}
    $\displaystyle \sum_{k=0}^\infty \sqrt{x} (x - 5)^k = \sqrt{x} +
    \sqrt{x} (x - 5) + \sqrt{x} (x - 5)^2 + \cdots$.  This is not a
    power series because of the $\sqrt{x}$.
    \end{solution}

    \item
    \begin{problem}[0.2in]
    $\displaystyle \sum_{k=5}^\infty \frac{(x - k)^k}{k}$
    \end{problem}

    \begin{solution}
    $\displaystyle \sum_{k=5}^\infty \frac{(x - k)^k}{k} = \frac{(x -
    5)^5}{5} + \frac{(x - 6)^6}{6} + \frac{(x - 7)^7}{7} + \cdots$.
    This is not a power series because there is no center.
    \end{solution}
  \end{enumerate}
  \probonly{}

\end{enumerate}


\begin{problemonly}
\begin{framed}
\textbf{\centerline{Theorem on the Convergence of a Power Series}} 
For a given power
series $\displaystyle \sum_{k=0}^\infty a_k (x - c)^k = a_0 + a_1 (x - c) +
  a_2 (x - c)^2 + \cdots $ centered at
$c$, there are 3 possibilities:
\begin{enumerate}[topsep=2pt,itemsep=0pt]
\item The series converges only when $x = c$.
\item The series converges absolutely for all values of $x$.
\item There is a positive number $R$ such that the series converges
absolutely when $|x - c| < R$ and diverges when $|x - c| > R$.  This distance $R$ from the center $c$
is called the \uline{radius of convergence}.  (Note that this doesn't
say anything about what happens when $|x - c| = R$.)
\end{enumerate}
If (1) is true, we say that the radius of convergence of the series is
0; if (2) is true, we say that the radius of convergence is $\infty$.
So, every power series has a non-negative radius of convergence.

\textbf{Terminology.}  The
\uline{largest open interval of convergence} is the open interval $(c
- R, c + R)$. The \uline{interval of convergence} of a power
series is the set of all $x$-values for which the power series converges.  
\end{framed}
\end{problemonly}

\begin{enumerate}[resume]
\pspace{\newpage}

\item 
 
 \begin{problem}[0.5in]
For each of the five Taylor series in problem 1, state the radius of convergence, the largest open interval of convergence, and the interval of convergence. 
\end{problem}

\begin{solution}
The Taylor series for $\frac{1}{1-x}$ and $\ln(1+x)$ both have a radius of convergence of 1, and a largest open interval of convergence of $(-1,1)$. The Taylor series for $e^x$, $\cos(x)$, and $\sin(x)$ all have a radius of convergence of $\infty$, and a largest open interval of convergence of $(-\infty, \infty)$.

The intervals of convergence for all 5 Taylor series are simply the answers to problems 1 and 2.
\end{solution}

\item 
\begin{enumerate}
  \item 
  \begin{problem}[2.5in]
  Find the radius of convergence of the power series $\displaystyle
  \sum_{k=1}^\infty \frac{(x - 3)^k}{k^2 \cdot 7^k}$.
  \end{problem}

  \begin{solution}
  When trying to find the radius of convergence or interval of
  convergence of a power series, we generally try the Ratio Test
  first.  The Ratio Test tells us to look at the following limit: 
  \begin{align*}
  \lim_{k \to \infty} \left| \frac{\frac{(x - 3)^{k + 1}}{(k + 1)^2
  \cdot 7^{k+1}}}{\frac{(x - 3)^k}{k^2 \cdot 7^k}} \right| &=
  \lim_{k \to \infty} \left| \frac{(x - 3) k^2}{(k + 1)^2 \cdot 7}
  \right| \\
  &= \left| \frac{x - 3}{7} \right| \lim_{k \to \infty} \left|
  \frac{k^2}{(k + 1)^2} \right| \\
  &= \left| \frac{x - 3}{7} \right|
  \end{align*}
  So, the Ratio Test says that the series converges when
  $\displaystyle \left| \frac{x - 3}{7} \right| < 1$ and diverges when
  $\displaystyle \left| \frac{x - 3}{7} \right| > 1$.  (When
  $\displaystyle \left| \frac{x - 3}{7} \right| = 1$, the Ratio Test
  gives no information.)  In other words, the series converges when
  $|x - 3| < 7$ and diverges when $|x - 3| > 7$.  This tells us that
  the \fbox{radius of convergence is $7$}.

  Here is a diagram showing what the Ratio Test tells us:
  \begin{center}
  \includegraphics[width=\linewidth]{ws/images/radius-interval-practice-4-ratio-test}
  \end{center}
  (Remember that we can read the inequality $|x - 3| < 7$ as saying,
  ``The distance between $x$ and $3$ is less than $7$.'')  The Ratio
  Test is inconclusive at the two blue points.
  \end{solution}

  \item
  \begin{problem}[0.25in]
  Find the largest open interval of convergence of the power series
  $\displaystyle \sum_{k=1}^\infty \frac{(x - 3)^k}{k^2 \cdot 7^k}$.
  \end{problem}

  \begin{solution}
  From the diagram, we can see that the largest open interval of
  convergence is the interval $\boxed{(-4, 10)}$.  (Notice that $3$,
  the center of the power series, is right in the middle of this
  interval; this is always the case.)
  \end{solution}

  \item
  \begin{problem}[2.5in]
  Find the interval of convergence of the power series 
  $\displaystyle \sum_{k=1}^\infty \frac{(x - 3)^k}{k^2 \cdot 7^k}$.
  \end{problem}

  \begin{solution}
  To finish finding the interval of convergence, we just need to
  figure out what happens at the two points $x = -4$ and $x = 10$.
  (These are the two points where the Ratio Test was inconclusive.)

  When $x = 10$, the given series is equal to $\displaystyle
  \sum_{k=1}^\infty \frac{(10 - 3)^k}{k^2 \cdot 7^k} = \sum_{k=1}^\infty
  \frac{7^k}{k^2 \cdot 7^k} = \sum_{k=1}^\infty \frac{1}{k^2}$, which
  converges, being a $p$-series with $p=2$.

  When $x = -4$, the given series is equal to $\displaystyle
  \sum_{k=1}^\infty \frac{(-4 - 3)^k}{k^2 \cdot 7^k} =
  \sum_{k=1}^\infty \frac{(-7)^k}{k^2 \cdot 7^k} = \sum_{k=1}^\infty
  \frac{(-1)^k}{k^2}$.  This series is absolutely convergent (because
  the series $\displaystyle \sum_{k=1}^\infty \left|
  \frac{(-1)^k}{k^2} \right|$ is just $\displaystyle
  \sum_{k=1}^\infty \frac{1}{k^2}$, which converges), and we know that
  every absolutely convergent series is also convergent.        

  So, the interval of convergence is $\boxed{[ -4, 10]}$. 
  \end{solution}

  \item
  \begin{problem}[0.35in]
  Draw the interval of convergence on a number line.
  \end{problem}

  \begin{solution}
  \begin{center}
  \includegraphics[width=5in]{ws/images/radius-interval-practice-4} 
  \end{center}
  Notice that the center of the power series (3) is in the center of
  the interval, as it should be.  The width of the interval (14) is
  twice the radius of convergence.
  
  Also, note that the interval of convergence for the power series is simply the domain of the function that the power series represents. So the power series in this question is a function whose domain is $[ -4, 10]$. 
  \end{solution}
\end{enumerate}


\pspace{\newpage}





\item 
\note{I've put this here just
so students have an example where the radius of convergence is 0.}

In this problem, you'll look at the power series $\displaystyle
\sum_{n=0}^\infty n! x^n$.

\begin{enumerate}
  \item
  \begin{problem}[1in]
  Write out the first few terms of the power series.  Does the series
  converge when $x = 0$?
  \end{problem}

  \begin{solution}
  The power series $\displaystyle \sum_{n=0}^\infty n! x^n$ is equal
  to $0! + 1! x + 2! x^2 + 3! x^3 + \cdots$, or $\boxed{1 + x + 2! x^2
  + 3!  x^3 + \cdots}$.  When $x = 0$, this is equal to $1 + 0 + 0 +
  \cdots$, which \fbox{converges}.
  \end{solution}

  \item
  \begin{problem}[3in]
  Find all values of $x$ for which the series $\displaystyle
  \sum_{n=0}^\infty n! x^n$ converges.  Make sure your answer is
  consistent with your answer to part (a).
  \end{problem}

  \begin{solution}
  Let's try the Ratio Test: $\displaystyle \lim_{n \to \infty} \left|
  \frac{(n + 1)! x^{n+1}}{n! x^n} \right| = |x| \lim_{n \to \infty}
  (n + 1)$.  If $x \neq 0$, then this limit is $\infty$, so the Ratio
  Test says that the series diverges.  If $x = 0$, then the limit is
  just $\displaystyle \lim_{n \to \infty} 0 = 0$, and the Ratio Test
  says that the series converges.  (This agrees with our answer to
  {{\#8}(a)}.)

  So, the series converges \fbox{only for $x = 0$}.
  \end{solution}

  \item
  \begin{problem}
  What is the radius of convergence of this power series?      
  \end{problem}

  \begin{solution}
  It is \fbox{0}, since the power series converges only at the center
  (that is, at $x = 0$).
  \end{solution}
\end{enumerate}

\pspace{\newpage}

\begin{problemonly}
\begin{framed}
\textbf{Definition.} If a function $f(x)$ is \textit{equal} to a power
series $\displaystyle \sum_{k=0}^\infty a_k (x - c)^k$ (possibly only
on a limited domain around $c$), we say that $\displaystyle
\sum_{k=0}^\infty a_k (x - c)^k$ is a \uline{power series representation
of $f(x)$ at $c$}.
\end{framed}
\end{problemonly}




\item 
\begin{enumerate}
  \item
  \begin{problem}[\vfill]
  Find a power series representation of the function $f(x) =
  \displaystyle \frac{1}{1 + 4x^2}$.  For what $x$ can we be sure that
  this representation is valid?
  \end{problem}

  \begin{solution}
  We know that
  \[ \frac{1}{1 - x} = 1 + x + x^2 + x^3 + \cdots = \sum_{k=0}^\infty
  x^k \textrm{ when } |x| < 1.\] If we replace $x$ by $-4x^2$
  everywhere in this statement, then we get
  \[ \displaystyle \frac{1}{1 + 4x^2} = 1 - 4x^2 + 16 x^4 - 64 x^6 +
  \cdots = \sum_{k=0}^\infty (-4x^2)^k \textrm{ when } |-4x^2| < 1.\]
  Equivalently,
  \begin{equation}
  \frac{1}{1 + 4x^2} = \sum_{k=0}^\infty (-4)^k x^{2k} \textrm{ when
  } |x| < \frac{1}{2}
  \end{equation}
  We can visualize the equation like this:
  \begin{center}
  \includegraphics[width=2in]{ws/images/manipulate-geometric-rep-1-eqn1} 
  \end{center}
  In this picture, the black curve is the graph of $\frac{1}{1 +
  4x^2}$; the blue one is the graph of the power series
  $\displaystyle \sum_{k=0}^\infty (-4)^k x^{2k}$.

  {\LARGE\Pointinghand} \textbf{Important Note:} When we talk about
  whether a series converges, we often write the series in sigma
  notation without specifying the starting index of summation.  This
  is fine, since we know that the beginning terms of a series have no
  effect on whether the series converges.  However, when we write
  power series representations, we \textbf{must} specify the starting
  index because we are talking about the sum of the series, not just
  whether the series converges.  So, for example, it is perfectly
  correct to write simply, ``$\sum x^k$ converges when $|x| < 1$.''
  However, if we want to talk about what the sum of the series is
  actually \textit{equal} to, we must specify the starting value of
  $k$.  For instance, $\displaystyle
  \sum_{k=\textcolor{red}{0}}^\infty x^k = \frac{1}{1 - x}$, but
  $\displaystyle \sum_{k=\textcolor{red}{1}}^\infty x^k = \frac{x}{1 -
  x}$.  Simply writing $\sum x^k = \frac{1}{1 - x}$ is incorrect.
  \end{solution}

  \item
  \note{Students are generally not sure what happens to the ``valid for
  $|x| < \frac{1}{2}$'' when we multiply by $x$.  Sometimes they try
  to multiply the inequality somehow to say that $|x^2| < \frac{1}{2}$
  or something.  I go back to the idea that the $|x| < \frac{1}{2}$ is
  telling us the domain of the power series, and that multiplying by
  $x$ does not affect the domain.}

  \begin{problem}[\vfill]
  Find a power series representation of the function $g(x) =
  \displaystyle \frac{x}{1 + 4x^2}$.  For what $x$ can we be sure that
  this representation is valid?
  \end{problem}

  \begin{solution}
  If we multiply both sides of (1)
  by $x$, we get that
  \[ \frac{x}{1 + 4x^2} = \sum_{k=0}^\infty (-4)^k x^{2k + 1} \textrm{
  when } |x| < \frac{1}{2}.\] You may wonder why this is still valid
  when $|x| < \frac{1}{2}$.  Remember that $|x| < \frac{1}{2}$ is
  telling us what the domain of the power series is; multiplying by
  $x$ does not change that domain.  It may help to think about this
  visually: if we multiply both graphs in the previous part by $x$,
  that doesn't change where they are defined, just their heights:
  \begin{center}
  \includegraphics[width=2in]{ws/images/manipulate-geometric-rep-1-eqn2} 
  \end{center}
  In this picture, the black curve is the graph of $\frac{x}{1 +
  4x^2}$; the blue one is the graph of the power series
  $\displaystyle \sum_{k=0}^\infty (-4)^k x^{2k+1}$. 
  \end{solution}

  \item 
  \begin{problem}[0.5in]
  In fact, the series you found in the two previous parts should have
  both been geometric.  Use your knowledge of geometric series to
  verify that your answers are correct.
  \end{problem}

  \begin{solution}
  In {{\#6}(a)}, we claimed that
  \[ \frac{1}{1 + 4x^2} = \sum_{k=0}^\infty (-4)^k x^{2k} \textrm{
  when } |x| < \frac{1}{2}.\] The series we found can be written in
  $\cdots$ notation as $1 - 4x^2 + 16x^4 - 64 x^6 + \cdots$; it's
  geometric with $a = 1$ and $r = -4x^2$, so it converges to
  $\frac{a}{1 - r} = \frac{1}{1 + 4x^2}$ when $|-4x^2| < 1$
  (equivalently, $|x| < \frac{1}{2}$).  Thus, our claim was correct.

  In {{\#6}(b)}, we claimed that
  \[ \frac{x}{1 + 4x^2} = \sum_{k=0}^\infty (-4)^k x^{2k + 1} \textrm{
  when } |x| < \frac{1}{2}.\] The series we found can be written in
  $\cdots$ notation as $x - 4x^3 + 16 x^5 - 64 x^7 + \cdots$; it's
  geometric with $a = x$ and $r = -4x^2$, so it converges to
  $\frac{a}{1 - r} = \frac{x}{1 + 4x^2}$ when $|-4x^2| < 1$
  (equivalently, $|x| < \frac{1}{2}$).  Thus, our claim was correct.
  \end{solution}
\end{enumerate}


\end{enumerate}




\end{document}